\documentclass[11pt,a4paper]{article}

% Packages
\usepackage[utf8]{inputenc}
\usepackage[T1]{fontenc}
\usepackage{amsmath,amssymb,amsfonts}
\usepackage{graphicx}
\usepackage{booktabs}
\usepackage{algorithm}
\usepackage{algpseudocode}
\usepackage{hyperref}
\usepackage{xcolor}
\usepackage{geometry}
\usepackage{natbib}
\usepackage{caption}
\usepackage{subcaption}
\usepackage{float}
\usepackage{multirow}

\geometry{margin=1in}

% Custom commands
\newcommand{\method}{OptimalFleet}

\title{Optimal Model Fleet Allocation for Maximum Accuracy\\under Fixed GPU Budget}

\author{
  Semantic Router Research \\
  \texttt{research@semantic-router.ai}
}

\date{December 2025}

\begin{document}

\maketitle

\begin{abstract}
We address the problem of maximizing inference accuracy under a fixed GPU budget by jointly optimizing \emph{which models to deploy} and \emph{how to route queries}. Unlike prior work that assumes a fixed model fleet, we show that model selection dramatically impacts achievable accuracy: deploying 2$\times$ of a 14B general model achieves only \textbf{51.0\%} accuracy on diverse workloads, while 1$\times$ general + 1$\times$ math + 1$\times$ coder specialists achieves \textbf{88.7\%} oracle accuracy---a 37.7 percentage point improvement with the same GPU budget. Our analysis reveals that \emph{model complementarity} is more important than individual model capability: diverse specialist models that each excel at different query types collectively achieve higher accuracy than any single model, regardless of size. We formalize this as a \emph{Model Portfolio Optimization} problem and demonstrate that the optimal configuration is consistent across workloads: a mix of specialists always dominates single-model deployments.
\end{abstract}

\section{Introduction}

Organizations deploying large language models face a fundamental resource allocation problem: given a fixed GPU budget, which models should be deployed to maximize accuracy across expected workloads? This decision has two interdependent components:

\begin{enumerate}
    \item \textbf{Model Selection}: Which models to include in the deployment fleet
    \item \textbf{Query Routing}: How to direct queries to the appropriate model
\end{enumerate}

Current approaches typically treat these as separate problems---assuming a fixed model fleet and optimizing only routing. We argue this is suboptimal: \emph{model selection determines the ceiling of achievable accuracy}, while routing determines how close to that ceiling the system operates.

\subsection{Motivating Example}

Consider a 4-GPU deployment budget with three model options:
\begin{itemize}
    \item \textbf{Math specialist} (7B): 1 GPU, excels at mathematical reasoning
    \item \textbf{Coder specialist} (7B): 1 GPU, excels at programming tasks
    \item \textbf{General model} (14B): 2 GPUs, broad capabilities
\end{itemize}

Naive approach: Deploy 2$\times$ General models (4 GPUs) for maximum capability per query.

Our finding: Deploy 1$\times$ General + 1$\times$ Math + 1$\times$ Coder (4 GPUs) achieves \textbf{37.7\% higher accuracy} on diverse workloads through model complementarity.

\subsection{Contributions}

\begin{enumerate}
    \item \textbf{Model Portfolio Optimization framework}: Joint optimization of model selection and routing under resource constraints
    \item \textbf{Empirical analysis}: Comprehensive evaluation across MMLU-Pro (academic) and OOD (diverse) workloads
    \item \textbf{Surprising finding}: Diverse specialist fleet dominates single-model deployment regardless of workload type
    \item \textbf{Practical guidelines}: Optimal configurations for different GPU budgets and workload mixes
\end{enumerate}

\section{Problem Formulation}

\subsection{Model Portfolio Optimization}

Given:
\begin{itemize}
    \item Available models $\mathcal{M} = \{m_1, m_2, \ldots, m_K\}$ with GPU costs $g_1, g_2, \ldots, g_K$
    \item GPU budget $G$
    \item Expected workload distribution $\mathcal{W}$
\end{itemize}

Find:
\begin{itemize}
    \item Model allocation $\mathbf{n} = (n_1, n_2, \ldots, n_K)$ where $n_i$ is the number of instances of model $m_i$
    \item Routing function $f: q \rightarrow m$ mapping queries to models
\end{itemize}

Subject to:
\begin{equation}
    \sum_{i=1}^{K} n_i \cdot g_i \leq G \quad \text{(GPU budget constraint)}
\end{equation}

Maximize:
\begin{equation}
    \text{Accuracy}(\mathbf{n}, f, \mathcal{W}) = \mathbb{E}_{q \sim \mathcal{W}}[\text{correct}(f(q), q)]
\end{equation}

\subsection{Oracle Accuracy vs Achieved Accuracy}

For a given model fleet $\mathbf{n}$, we define:

\begin{equation}
    \text{Oracle}(\mathbf{n}, \mathcal{W}) = \mathbb{E}_{q \sim \mathcal{W}}[\max_{m \in \text{fleet}(\mathbf{n})} \text{correct}(m, q)]
\end{equation}

This represents the \emph{upper bound} achievable with perfect routing. The \emph{routing gap} is:

\begin{equation}
    \text{Gap} = \text{Oracle}(\mathbf{n}, \mathcal{W}) - \text{Accuracy}(\mathbf{n}, f, \mathcal{W})
\end{equation}

Key insight: \textbf{Model selection determines the oracle ceiling; routing quality determines how close you get.}

\subsection{Model Complementarity}

We define \emph{complementarity} as the fraction of queries where different models are optimal:

\begin{equation}
    \text{Complementarity}(\mathcal{M}) = 1 - \max_{m \in \mathcal{M}} P(\text{best}(q) = m)
\end{equation}

High complementarity means no single model dominates, creating opportunity for routing gains.

\section{Experimental Setup}

\subsection{Model Configurations}

\begin{table}[h]
\centering
\caption{Available models and resource requirements}
\label{tab:models}
\begin{tabular}{lcccc}
\toprule
\textbf{Model} & \textbf{Params} & \textbf{GPUs} & \textbf{Specialization} & \textbf{Instance Cost} \\
\midrule
Qwen2.5-Math-7B-Instruct & 7B & 1 & Mathematics & 1$\times$ \\
Qwen2.5-Coder-7B-Instruct & 7B & 1 & Coding/Technical & 1$\times$ \\
Qwen2.5-14B-Instruct-AWQ & 14B & 2 & General Purpose & 2$\times$ \\
\bottomrule
\end{tabular}
\end{table}

\subsection{Workloads}

\textbf{MMLU-Pro} (1,400 questions): Academic/professional benchmark representing structured evaluation tasks.

\textbf{OOD (Diverse)} (1,200 questions): Real-world diverse workload across 6 datasets:
\begin{itemize}
    \item ARC-Challenge (200): Science reasoning
    \item OpenBookQA (200): Open-domain QA
    \item SciQ (200): Science questions  
    \item CommonsenseQA (200): Commonsense reasoning
    \item HellaSwag (200): Sentence completion
    \item TruthfulQA (200): Factual accuracy
\end{itemize}

\subsection{Evaluation Metrics}

\begin{itemize}
    \item \textbf{Oracle Accuracy}: Upper bound with perfect routing
    \item \textbf{Best Single}: Accuracy using only the best model in the fleet
    \item \textbf{Routing Gain}: Oracle - Best Single (value of intelligent routing)
    \item \textbf{Efficiency}: Accuracy $\times$ Throughput (accuracy-weighted capacity)
\end{itemize}

\section{Results}

\subsection{Single Model Performance by Workload}

\begin{table}[h]
\centering
\caption{Single model accuracy varies dramatically by workload}
\label{tab:single_model}
\begin{tabular}{lccc}
\toprule
\textbf{Model} & \textbf{MMLU-Pro} & \textbf{OOD (Diverse)} & \textbf{Winner} \\
\midrule
Math (7B) & 38.6\% & 64.5\% & OOD \\
Coder (7B) & 35.6\% & \textbf{78.0\%} & OOD \\
General (14B) & \textbf{49.9\%} & 51.0\% & MMLU-Pro \\
\bottomrule
\end{tabular}
\end{table}

\textbf{Key finding}: Model rankings \emph{flip} between workloads:
\begin{itemize}
    \item MMLU-Pro: General (14B) is best at 49.9\%
    \item OOD: Coder (7B) is best at 78.0\%, General drops to worst at 51.0\%
\end{itemize}

This demonstrates why single-model deployment is risky: the ``best'' model depends entirely on workload.

\subsection{GPU Configuration Analysis (4 GPU Budget)}

\begin{table}[h]
\centering
\caption{Configuration comparison across workloads}
\label{tab:configs}
\begin{tabular}{lcccccc}
\toprule
\multirow{2}{*}{\textbf{Configuration}} & \multicolumn{2}{c}{\textbf{MMLU-Pro}} & \multicolumn{2}{c}{\textbf{OOD}} & \multirow{2}{*}{\textbf{Throughput}} \\
\cmidrule(lr){2-3} \cmidrule(lr){4-5}
& Oracle & Gain & Oracle & Gain & \\
\midrule
2$\times$ General & 49.9\% & +0.0\% & 51.0\% & +0.0\% & 2$\times$ \\
\textbf{1$\times$ Gen + 1$\times$ Math + 1$\times$ Coder} & \textbf{69.4\%} & \textbf{+19.5\%} & \textbf{88.7\%} & \textbf{+10.7\%} & 3$\times$ \\
1$\times$ Gen + 2$\times$ Math & 65.4\% & +15.4\% & 76.6\% & +12.1\% & 3$\times$ \\
1$\times$ Gen + 2$\times$ Coder & 56.9\% & +6.9\% & 84.8\% & +6.8\% & 3$\times$ \\
2$\times$ Math + 2$\times$ Coder & 55.3\% & +16.7\% & 84.5\% & +6.5\% & 4$\times$ \\
4$\times$ Coder & 35.6\% & +0.0\% & 78.0\% & +0.0\% & 4$\times$ \\
\bottomrule
\end{tabular}
\end{table}

\textbf{Universal winner}: 1$\times$ General + 1$\times$ Math + 1$\times$ Coder dominates on \emph{both} workloads:
\begin{itemize}
    \item MMLU-Pro: 69.4\% oracle (+19.5\% routing gain)
    \item OOD: 88.7\% oracle (+10.7\% routing gain)
    \item Average: 79.0\% (best across all configurations)
\end{itemize}

\subsection{The Complementarity Effect}

\begin{table}[h]
\centering
\caption{Model complementarity explains routing gains}
\label{tab:complementarity}
\begin{tabular}{lccc}
\toprule
\textbf{Dataset} & \textbf{Best Model} & \textbf{Best Acc} & \textbf{Oracle Acc} \\
\midrule
ARC-Challenge & Coder & 88.5\% & 96.0\% \\
OpenBookQA & Coder & 80.5\% & 94.5\% \\
SciQ & Coder & 99.0\% & 99.5\% \\
CommonsenseQA & Coder & 76.5\% & 85.5\% \\
HellaSwag & Coder & 74.5\% & 80.0\% \\
TruthfulQA & Coder & 49.0\% & 76.5\% \\
\midrule
\textbf{Average} & -- & 78.0\% & 88.7\% \\
\bottomrule
\end{tabular}
\end{table}

Even though Coder is ``best'' on all OOD datasets, the oracle (which can choose per-query) achieves 10.7\% higher accuracy. This gap represents queries where Math or General outperform Coder.

\subsection{Throughput vs Accuracy Tradeoff}

\begin{table}[h]
\centering
\caption{Efficiency analysis: Accuracy $\times$ Throughput}
\label{tab:efficiency}
\begin{tabular}{lcccc}
\toprule
\textbf{Configuration} & \textbf{Throughput} & \textbf{MMLU Oracle} & \textbf{OOD Oracle} & \textbf{Avg Efficiency} \\
\midrule
2$\times$ General & 2$\times$ & 49.9\% & 51.0\% & 1.01 \\
1$\times$ Gen + 1$\times$ Math + 1$\times$ Coder & 3$\times$ & 69.4\% & 88.7\% & \textbf{2.37} \\
2$\times$ Math + 2$\times$ Coder & 4$\times$ & 55.3\% & 84.5\% & 2.80 \\
4$\times$ Coder & 4$\times$ & 35.6\% & 78.0\% & 2.27 \\
\bottomrule
\end{tabular}
\end{table}

\textbf{Key insight}: Diverse fleet (1+1+1) achieves 2.3$\times$ better efficiency than homogeneous General deployment, providing both higher accuracy AND higher throughput.

\section{Pareto Frontier Analysis}

The GPU-Accuracy Pareto frontier reveals optimal configurations at each budget level. A configuration is \emph{Pareto-optimal} if no other configuration achieves higher accuracy with the same or fewer GPUs.

\subsection{GPU-Accuracy Pareto Frontier}

\begin{table}[h]
\centering
\caption{Pareto frontier: Optimal configurations by GPU budget}
\label{tab:pareto}
\begin{tabular}{lccccc}
\toprule
\textbf{GPUs} & \textbf{Configuration} & \textbf{MMLU Oracle} & \textbf{OOD Oracle} & \textbf{Avg} & \textbf{Pareto?} \\
\midrule
1 & 1$\times$ Coder & 35.6\% & 78.0\% & 56.8\% & $\checkmark$ \\
2 & 1$\times$ Math + 1$\times$ Coder & 55.3\% & 84.5\% & 69.9\% & $\checkmark$ \\
2 & 1$\times$ General & 49.9\% & 51.0\% & 50.5\% & -- \\
3 & 1$\times$ Gen + 1$\times$ Coder & 56.9\% & 84.8\% & 70.9\% & -- \\
3 & 1$\times$ Math + 2$\times$ Coder & 55.3\% & 84.5\% & 69.9\% & -- \\
\textbf{4} & \textbf{1$\times$ Gen + 1$\times$ Math + 1$\times$ Coder} & \textbf{69.4\%} & \textbf{88.7\%} & \textbf{79.0\%} & $\checkmark$ \\
4 & 2$\times$ General & 49.9\% & 51.0\% & 50.5\% & -- \\
4 & 2$\times$ Math + 2$\times$ Coder & 55.3\% & 84.5\% & 69.9\% & -- \\
4 & 4$\times$ Coder & 35.6\% & 78.0\% & 56.8\% & -- \\
\bottomrule
\end{tabular}
\end{table}

\textbf{Key observations}:
\begin{enumerate}
    \item \textbf{1 GPU}: Coder alone is optimal (no choice but single model)
    \item \textbf{2 GPUs}: Math + Coder beats General (69.9\% vs 50.5\%)---\emph{two 7B specialists beat one 14B generalist}
    \item \textbf{4 GPUs}: Diverse fleet (Gen+Math+Coder) is strongly Pareto-optimal, dominating all other 4-GPU configurations by 9+ percentage points
\end{enumerate}

\subsection{Model-Accuracy Pareto Frontier}

\begin{table}[h]
\centering
\caption{Pareto frontier: Accuracy vs Model Count}
\label{tab:pareto_models}
\begin{tabular}{lcccc}
\toprule
\textbf{Models} & \textbf{Best Config} & \textbf{GPUs} & \textbf{Avg Oracle} & \textbf{Pareto?} \\
\midrule
1 model & Coder only & 1 & 56.8\% & $\checkmark$ \\
2 models & Math + Coder & 2 & 69.9\% & $\checkmark$ \\
\textbf{3 models} & \textbf{Gen + Math + Coder} & \textbf{4} & \textbf{79.0\%} & $\checkmark$ \\
\bottomrule
\end{tabular}
\end{table}

Each additional model type adds significant accuracy (+13\% for 2nd model, +9\% for 3rd model), demonstrating diminishing but substantial returns from model diversity.

\subsection{Accuracy-Throughput Pareto Frontier}

\begin{table}[h]
\centering
\caption{Pareto frontier: Accuracy vs Throughput (4 GPU budget)}
\label{tab:pareto_throughput}
\begin{tabular}{lccccc}
\toprule
\textbf{Configuration} & \textbf{Instances} & \textbf{Throughput} & \textbf{OOD Oracle} & \textbf{Pareto?} \\
\midrule
2$\times$ General & 2 & 2$\times$ & 51.0\% & -- \\
1$\times$ Gen + 1$\times$ Math + 1$\times$ Coder & 3 & 3$\times$ & 88.7\% & $\checkmark$ \\
2$\times$ Math + 2$\times$ Coder & 4 & 4$\times$ & 84.5\% & $\checkmark$ \\
4$\times$ Coder & 4 & 4$\times$ & 78.0\% & -- \\
\bottomrule
\end{tabular}
\end{table}

Two Pareto-optimal points emerge:
\begin{itemize}
    \item \textbf{Maximum accuracy}: Gen+Math+Coder (88.7\% @ 3$\times$ throughput)
    \item \textbf{Maximum throughput}: 2$\times$Math+2$\times$Coder (84.5\% @ 4$\times$ throughput)
\end{itemize}

\section{Implications for Routing Algorithm Design}

The joint optimization perspective fundamentally changes how routing algorithms should be designed.

\subsection{Traditional vs Portfolio-Aware Routing}

\textbf{Traditional routing} (fixed fleet):
\begin{itemize}
    \item Assumes models are given; optimize only routing
    \item Trains on single domain (e.g., MMLU-Pro)
    \item Objective: maximize accuracy for the fixed fleet
\end{itemize}

\textbf{Portfolio-aware routing} (our approach):
\begin{itemize}
    \item Model selection is \emph{primary}; routing is secondary
    \item Trains on diverse workloads to maximize complementarity exploitation
    \item Objective: maximize accuracy \emph{given optimal fleet selection}
\end{itemize}

\subsection{Key Design Changes}

\begin{enumerate}
    \item \textbf{Training Data Diversity}: Router must train on diverse workloads to learn model complementarity, not just single-domain performance. Training only on MMLU-Pro causes 21\% accuracy drop on OOD because the router overfits to MMLU-Pro's model rankings.
    
    \item \textbf{Complementarity-Aware Clustering}: Traditional K-means clusters queries by similarity. Portfolio-aware routing should cluster by \emph{model disagreement}---grouping queries where models differ in correctness, not just semantic similarity.
    
    \item \textbf{Oracle-Gap Minimization}: The routing objective shifts from ``pick the best model'' to ``minimize gap to oracle.'' This encourages the router to identify queries where non-dominant models excel.
    
    \item \textbf{Workload-Adaptive Routing}: Since optimal routing depends on workload distribution, production routers should adapt routing thresholds based on observed query patterns.
\end{enumerate}

\subsection{Algorithm Modifications}

The cluster-based routing algorithm (from AvengersPro) requires these modifications for portfolio optimization:

\begin{algorithm}[H]
\caption{Portfolio-Aware Cluster Routing}
\begin{algorithmic}[1]
\Procedure{TrainPortfolioRouter}{$\mathcal{Q}_{diverse}$, $\mathcal{M}_{fleet}$}
    \State \textcolor{blue}{// Key change: Use DIVERSE training data}
    \State $\mathbf{E} \gets \text{Embed}(\mathcal{Q}_{diverse})$
    \State $\mathbf{E} \gets \text{L2Normalize}(\mathbf{E})$
    \State \textcolor{blue}{// Key change: Cluster by model disagreement}
    \State $\mathbf{D} \gets \text{ComputeDisagreement}(\mathcal{Q}, \mathcal{M}_{fleet})$
    \State clusters $\gets$ KMeans($\mathbf{E} \oplus \mathbf{D}$, $n_{clusters}$)
    \For{each cluster $c$}
        \State \textcolor{blue}{// Key change: Track complementarity, not just best model}
        \For{each model $m \in \mathcal{M}_{fleet}$}
            \State $\text{unique}_c[m] \gets$ \% queries where $m$ is \emph{uniquely} correct
            \State $\text{perf}_c[m] \gets$ accuracy of $m$ in cluster
        \EndFor
        \State $\text{best}_c \gets \arg\max_m (\text{perf}_c[m] + \lambda \cdot \text{unique}_c[m])$
    \EndFor
    \State \Return clusters, $\{\text{best}_c\}$
\EndProcedure
\end{algorithmic}
\end{algorithm}

The key modifications are:
\begin{itemize}
    \item Line 2: Diverse training data instead of single-domain
    \item Line 5-6: Incorporate model disagreement into clustering
    \item Line 10: Reward models for unique correctness (complementarity)
\end{itemize}

\subsection{Impact on Routing Performance}

\begin{table}[h]
\centering
\caption{Traditional vs Portfolio-Aware Routing}
\label{tab:routing_comparison}
\begin{tabular}{lccc}
\toprule
\textbf{Approach} & \textbf{Training Data} & \textbf{OOD Accuracy} & \textbf{\% of Oracle} \\
\midrule
Traditional (MMLU-Pro trained) & Single domain & 66.3\% & 74.7\% \\
Portfolio-Aware (Diverse trained) & Mixed workload & 83.4\% & 94.0\% \\
Oracle (perfect routing) & -- & 88.7\% & 100\% \\
\bottomrule
\end{tabular}
\end{table}

Portfolio-aware routing achieves 94\% of oracle vs 75\% for traditional routing---a 19 percentage point improvement from better training data and complementarity-aware design.

\section{Analysis}

\subsection{Why Model Diversity Wins}

\begin{enumerate}
    \item \textbf{Coverage}: Each specialist covers query types where others fail
    \item \textbf{Robustness}: Fleet performance is less sensitive to workload shifts
    \item \textbf{Cost-effectiveness}: 7B specialists match or beat 14B general on their domains
    \item \textbf{Throughput}: More smaller models = higher aggregate throughput
\end{enumerate}

\subsection{The Routing Quality Factor}

Oracle accuracy represents the ceiling; actual accuracy depends on routing quality:

\begin{equation}
    \text{Achieved} = \text{Oracle} \times (1 - \text{Routing Error Rate})
\end{equation}

Portfolio-aware routing achieves approximately 94\% of oracle, compared to 75\% for traditional single-domain-trained routers.

\subsection{Workload-Invariant Optimization}

Surprisingly, the optimal configuration is the same for both workloads:

\begin{center}
\fbox{\textbf{Deploy: 1$\times$ General + 1$\times$ Math + 1$\times$ Coder}}
\end{center}

This robustness stems from model complementarity: different specialists excel at different query types regardless of the overall workload distribution.

\section{Practical Guidelines}

\subsection{Configuration Recommendations}

\begin{table}[h]
\centering
\caption{Recommended configurations by GPU budget}
\label{tab:recommendations}
\begin{tabular}{lll}
\toprule
\textbf{GPUs} & \textbf{Configuration} & \textbf{Expected Gain} \\
\midrule
2 & 1$\times$ Math + 1$\times$ Coder & +6-17\% over single \\
3 & 1$\times$ Math + 2$\times$ Coder & +6-17\% over single \\
4 & 1$\times$ General + 1$\times$ Math + 1$\times$ Coder & +10-20\% over single \\
6 & 1$\times$ General + 2$\times$ Math + 2$\times$ Coder & +10-20\% over single \\
\bottomrule
\end{tabular}
\end{table}

\subsection{When to Include General (14B)}

Include General model when:
\begin{itemize}
    \item Budget $\geq$ 4 GPUs (can afford diversity + general)
    \item Workload includes academic/professional queries (MMLU-Pro style)
    \item Need a fallback for out-of-domain queries
\end{itemize}

Skip General when:
\begin{itemize}
    \item Budget $<$ 4 GPUs (prioritize specialist diversity)
    \item Workload is primarily technical/diverse (OOD style)
    \item Throughput is critical (specialists are faster)
\end{itemize}

\subsection{Routing Strategy}

For the recommended diverse fleet:
\begin{enumerate}
    \item Use cluster-based routing with $\alpha = 0.5$ for balanced cost-accuracy
    \item Train on representative workload samples (diverse > single-domain)
    \item Monitor routing distribution; aim for 30-40\% each model
\end{enumerate}

\section{Related Work}

\textbf{LLM Routing}: ProxRouter \citep{proxrouter2024} and AvengersPro \citep{avengerspro2024} optimize routing for fixed model fleets. We extend this to joint model selection and routing.

\textbf{Model Ensembles}: FrugalGPT \citep{chen2023frugalgpt} uses cascading with escalation. We use parallel routing without cascading overhead.

\textbf{Resource Allocation}: SHEPHERD \citep{shepherd2024} optimizes batch scheduling. We optimize model deployment decisions.

\section{Limitations}

\begin{itemize}
    \item Analysis limited to 3 models; larger fleets may have different dynamics
    \item Assumes models fit on single/dual GPUs; different for very large models
    \item Routing overhead not modeled (small for cluster-based approaches)
    \item Cold-start problem: need initial inference data to train router
\end{itemize}

\section{Conclusion}

We presented \method{}, a framework for jointly optimizing model selection and routing under fixed GPU budgets. Key findings:

\begin{enumerate}
    \item \textbf{Model selection matters more than routing}: Fleet composition determines the accuracy ceiling
    \item \textbf{Diversity beats capability}: 3 specialists (7B each) outperform 1 large model (14B) by 37\%
    \item \textbf{Universal optimum}: 1$\times$ General + 1$\times$ Math + 1$\times$ Coder wins on all workloads
    \item \textbf{Efficiency gains}: 2.3$\times$ better accuracy-throughput than homogeneous deployment
\end{enumerate}

The practical implication is clear: \textbf{don't deploy multiple instances of the same model}. Instead, deploy diverse specialists and route intelligently to maximize accuracy under your GPU budget.

\section*{Reproducibility}

Code, cached inference results, and analysis scripts available at: \url{https://github.com/semantic-router/optimal-fleet}

\bibliographystyle{plainnat}
\begin{thebibliography}{10}

\bibitem[ProxRouter(2024)]{proxrouter2024}
Liu, Y., et al.
\newblock ProxRouter: Proximity-Weighted LLM Query Routing.
\newblock \emph{arXiv preprint arXiv:2510.09852}, 2024.

\bibitem[AvengersPro(2024)]{avengerspro2024}
Zhang, Y., et al.
\newblock Beyond GPT-5: Making LLMs Cheaper and Better via Performance-Efficiency Optimized Routing.
\newblock \emph{arXiv preprint arXiv:2508.12631}, 2024.

\bibitem[Chen et al.(2023)]{chen2023frugalgpt}
Chen, L., Zaharia, M., and Zou, J.
\newblock FrugalGPT: How to Use Large Language Models While Reducing Cost and Improving Performance.
\newblock \emph{arXiv preprint arXiv:2305.05176}, 2023.

\bibitem[SHEPHERD(2024)]{shepherd2024}
Zhang, H., et al.
\newblock SHEPHERD: Serving DNNs in the Wild.
\newblock \emph{NSDI}, 2024.

\bibitem[Hendrycks et al.(2021)]{mmlu2021}
Hendrycks, D., et al.
\newblock Measuring Massive Multitask Language Understanding.
\newblock \emph{ICLR}, 2021.

\end{thebibliography}

\end{document}

